%!TeX root=20_Hauptteil.tex

%\chapter{Implementierung der Positionierung}
%
%\subsection{Kommunikation PX4}
%
%\subsection{Flugsteuerung}
%
%\section{Bildaufnahme}
%
%\subsection{Kamerakalibrierung}
%
%\section{Entzerrung}
%
%\section{Bildverarbeitung}
%
%\section{Positionsbestimmung}
%
%\subsection{extrinsische Kamerakalibrierung}

%!TeX root=20_Hauptteil.tex
\chapter{Implementierung der Positionierung}
\label{cap:lokalisierung}
Dieses Kapitel beschreibt die im Flug ablaufende Positionsbestimmung des Fluggeräts.
Im Gegensatz zur Kartierung (Kapitel~\ref{cap:kartierung}) läuft die Verarbeitung hier zur Laufzeit auf dem Companion-Computer.
Die Kette folgt dem in Abschnitt~\ref{sec:softwarearchitektur} dargestellten Datenfluss: Bildaufnahme, Entzerrung, Bildverarbeitung, Positionsbestimmung, 
Kommunikation mit der Flugsteuerung und schließlich die Flugsteuerung selbst.
Die ersten drei Schritte entsprechen methodisch denen der Kartierung, werden aber mit höherer Bildrate und auf der GPU des Jetson ausgeführt.
Die gesamte Positionsbestimmung erfolgt im \gls{enu}/\gls{flu}-System, in dem auch die Karte vorliegt, und wird erst am 
Ende der Verarbeitungskette in das von PX4 benötigte \gls{ned}/\gls{frd}-System umgerechnet.



\section{Softwarearchitektur der Positionierung}
\label{sec:softwarearchitektur_positionierung}
Die Positionierung ist auf zwei Ausführungsumgebungen aufgeteilt: den Host des Companion-Computers und einen Container.
Diese Trennung ergibt sich aus der Verwendung von Isaac~ROS für die hardwarebeschleunigten Schritte.
Diese Beschleunigung ist notwendig, um der geforderten Übertragungsrate des \gls{px4} gerecht zu werden. 
Isaac~ROS ist eine von NVIDIA bereitgestellte Sammlung von ROS2-Paketen, die für die Ausführung auf den \glspl{gpu} der Jetson-Plattform optimiert sind.
Die Pakete werden in einem vorkonfigurierten Container ausgeführt, der die benötigten Bibliotheken und Treiber bereitstellt.
Innerhalb dieses Containers werden die Entzerrung und die Markererkennung auf der \gls{gpu} ausgeführt, was gegenüber 
einer CPU-basierten Verarbeitung eine höhere Bildrate ermöglicht.
Die übrigen Schritte laufen direkt auf dem Host des Jetson.
Dort erfolgen die Bildaufnahme über gscam, die Positionsbestimmung aus den Markerdetektionen sowie die Kommunikation mit der Flugsteuerung.
Vor der Eigentlichen Positionierung wird eine extrinische Kalibrierung der verwendeten Kameras durchgeführt.
Container und Host kommunizieren über ROS2-Topics.
Die Bildaufnahme stellt das Kamerabild auf \topic{/camera_1/image_raw} mit den zugehörigen Parametern auf \topic{/camera_1/camera_info} bereit.
Im Container entzerrt der Knoten \texttt{RectifyNode} das Bild und gibt es als \topic{/camera_1/image_rect} aus.
Der Knoten \texttt{AprilTagNode} detektiert darauf die Marker und veröffentlicht die Detektionen auf \topic{/camera_1/tag_detections}.
Auf dem Host bestimmt der Knoten \texttt{vision\_pose\_uxrce\_node} aus diesen Detektionen die Pose des Fluggeräts und übergibt sie 
über \topic{/fmu/in/vehicle_visual_odometry} an die Flugsteuerung.
Abbildung~ref{fig:architektur_positionierung} zeigt die Architektur der Implementierung zur Positionierung.
%den oberen Satz nochmal mit den Bild optimieren.....

% TODO: Abbildung -- Knoten und Topics, getrennt nach Host und Container.
\begin{figure}[htbp]
    \centering
    % \includegraphics[width=\textwidth]{Bilder/Architektur_Positionierung.png}
    \caption{Knoten und Topics der Positionierung, aufgeteilt nach Host und Container}
    \label{fig:architektur_positionierung}
\end{figure}
\FloatBarrier





\section{Bildaufnahme Positionierung}
\label{sec:pos_bildaufnahme}
Die Bildaufnahme erfolgt wie bei der Kartierung (Abschnitt~\ref{sec:kartierung_bildaufnahme}) über einen gscam-Knoten.
Auch hier wird wie bei der Implementierung der Kartierung exemplarisch die Logitech~C922 verwendet.
Wie in Abschnitt~\ref{Kamera_Kalibrierung} beschrieben, muss die Kalibrierung für jede Kamera und für die hier verwendete 
Auflösung von $1280 \times 720$ Pixeln durchgeführt werden.


Listing~\ref{lst:gscam_node_pos} zeigt die geänderte Konfiguration.

\begin{lstlisting}[language=Python, caption={Parameter des gscam-Knotens zur Bildaufnahme im Flug}, label={lst:gscam_node_pos}]
'gscam_config': (
    'v4l2src device=/dev/video0 do-timestamp=true ! '
    'image/jpeg,width=1280,height=720,framerate=60/1 ! '
    'queue leaky=downstream max-size-buffers=1 ! '
    'nvv4l2decoder mjpeg=1 ! '
    'queue leaky=downstream max-size-buffers=1 ! '
    'nvvidconv ! '
    'video/x-raw,format=BGRx ! '
    'videoconvert ! '
    'video/x-raw,format=RGB,framerate=60/1'
),
'camera_name': 'camera_1',
'frame_id': 'camera_1_optical_frame',
'camera_info_url': camera_info_url,
'use_gst_timestamps': True,
'sync_sink': False,
'preroll': False,
'image_encoding': 'rgb8',
'use_sensor_data_qos': False,
\end{lstlisting}

Die beiden wesentlichen Unterschiede betreffen die Bildrate und die Dekodierung.
Die Kamera wird hier mit einer Auflösung von $1280 \times 720$ Pixeln bei $60$ \gls{fps} betrieben.
Die deutlich höhere Bildrate ist notwendig, da die Positionsbestimmung im Flug fortlaufend erfolgt und eine geringe Latenz erfordert.
Die Auflösung von $1280 \times 720$ Pixeln wurde gewählt, da die Logitech~C922 nur bei dieser Auflösung $60$ \gls{fps} liefert.
Möglich wird die hohe Bildrate durch die hardwarebeschleunigte Dekodierung des MJPEG-Streams über \texttt{nvv4l2decoder} und \texttt{nvvidconv} 
auf der \gls{gpu} des Jetson.
Ohne die hardwarebeschleunigte Dekodierung sinkt die Bildrate über die Zeit, da die CPU mit der Dekodierung 
nicht hinterherkommt und ein Puffer in der Pipeline volläuft.


\section{Entzerrung}
\label{sec:pos_entzerrung}
Die Entzerrung erfolgt nach demselben Prinzip wie bei der Kartierung (Abschnitt~\ref{sec:kartierung_entzerrung}).
Im Unterschied dazu wird sie jedoch nicht über \texttt{rectify\_node} aus dem Paket image\_proc, sondern über den Knoten \texttt{RectifyNode} aus 
dem Paket isaac\_ros\_image\_proc umgesetzt.
Dieser läuft \gls{gpu}-beschleunigt innerhalb des Isaac-Containers.

Der Knoten liest das Bild von \topic{/camera_1/image_raw} mit den zugehörigen Parametern von \topic{/camera_1/camera_info} ein und gibt das 
entzerrte Bild als \topic{/camera_1/image_rect} mit den Parametern \topic{/camera_1/camera_info_rect} aus.
Wie bei der Kartierung entsprechen die intrinsischen Parameter des entzerrten Bildes der Projektionsmatrix aus 
Gleichung~\eqref{eq:Projektionsmatrix}, und die Verzeichnungskoeffizienten sind null.



\section{Detektion von AprilTags}
\label{sec:pos_detektion_apriltags}
Im Gegensatz zur Kartierung (Abschnitt~\ref{sec:offline_detektion_apriltags}) erfolgt die Detektion hier nicht offline, sondern fortlaufend während des Flugs.
Verwendet wird daher nicht \texttt{sync\_and\_detect\_node} aus dem Paket TagSLAM, sondern der Knoten \texttt{AprilTagNode} aus dem Paket isaac\_ros\_apriltag.
Dieser läuft ebenfalls \gls{gpu}-beschleunigt innerhalb des Isaac-Containers und ermöglicht so die Detektion bei voller Bildrate.
Der Knoten liest das entzerrte Bild von \topic{/camera_1/image_rect} mit den Parametern von \topic{/camera_1/camera_info_rect} ein und detektiert
darauf die AprilTag-Marker.
Über die Parameter werden die Markerfamilie \texttt{tag36h11}, die physische Markergröße von $0{,}167$ Metern sowie die maximale Anzahl von $40$ Markern
pro Bild festgelegt. Der hardwarebeschleunigte Detektor ist in der Lage nur eine einzige Markergröße zu detektieren, weshalb in Abschnitt
\ref{sec:pos_positionsbestimmung} eine Umrechnung anderer Größen erfolgt.
Für jeden erkannten Marker liefert der Knoten dessen Identifikationsnummer, den Mittelpunkt sowie die Bildkoordinaten der vier Eckpunkte.
Aus den Eckpunkten und der bekannten Markergröße schätzt der Detektor zusätzlich die relative Pose des Markers im Kamerakoordinatensystem.
Marker, die im Bild groß dargestellt werden, lassen sich dabei zuverlässiger detektieren als kleine.
Diese Pose setzt sich aus einem Translationsvektor und einer als Quaternion kodierten
Orientierung zusammen. Der Node gibt (abweichend von dem in \gls{ros2} Standardpaket) keine Kovarianzmatrix aus, die angibt, wie vertrauenswürdig die Schätzung ist.
Alle genannten Größen werden gebündelt als Nachricht vom Typ \texttt{AprilTagDetectionArray} auf dem Topic \topic{/camera_1/tag_detections} ausgegeben.
Die so bestimmte relative Markerpose bildet die Eingangsgröße für die Positionsbestimmung in Abschnitt~\ref{sec:pos_positionsbestimmung}.



\section{Extrinsische Kalibrierung zweier Kameras auf den PX4-Mittelpunkt}
\label{subsec:extrinsische_kalibrierung}
Die extrinsische Kalibrierung bestimmt für beide Kameras die statische Transformation vom \gls{px4}-Mittelpunkt zur Kamera, 
$^{\mathrm{D}}\mathbf{T}_{\mathrm{C}}$.
Diese Transformationsmatrizen werden zur Laufzeit von der Positionsbestimmung (Abschnitt~\ref{sec:pos_positionsbestimmung}) benötigt.
Die Kalibrierung findet zur Reproduzierbarkeit in einer Kalibrierbox statt.
Der Multikopter wird in dieser Box an einer fest vorgegebenen Position platziert und anhand von Wandmarkierungen ausgerichtet 
(Abbildung~\ref{fig:kalibrierbox}).
In der Box befinden sich zwei Marker, die zueinander sowie zum \gls{px4}-Mittelpunkt von Hand vermessen wurden.
Je eine Kamera ist auf je einen der Marker gerichtet.
Die drei vermessenen Positionen werden dem Node über eine Datei übergeben.
Zum Starten der Kalibrierung müssen zusätzlich \texttt{gscam} und \texttt{rectify} für beide Kameras gestartet werden.
%\begin{figure}[htbp]
%\centering
%\includegraphics[width=1.0\textwidth]{Bilder/Kalibrierbox.png}
%\caption{Drohne in der Kalibrierbox mit von Hand vermessenen Markern und den Wandmarkierungen zur Ausrichtung}
%\label{fig:kalibrierbox}
%\end{figure}
%\FloatBarrier
Sämtliche Posen werden im Koordinatensystem~$\mathrm{W}$ der Kalibrierbox angegeben.
Die Pose des \gls{px4}-Mittelpunkts in diesem Koordinatensystem ist $^{\mathrm{W}}\mathbf{T}_{\mathrm{D}}$.
Die in Abschnitt~\ref{subsec:pos_posenberechnung} beschriebene Posenberechnung liefert die Pose der jeweiligen Kamera in diesem Koordinatensystem, $^{\mathrm{W}}\mathbf{T}_{\mathrm{C}}$.
Die gesuchte Transformation ergibt sich durch Verkettung der Kamerapose mit der invertierten Pose des \gls{px4}-Mittelpunkts gemäß den Gleichungen~\eqref{eq:trafo_inverse} und~\eqref{eq:trafo_verkettung}:
\begin{equation}
^{\mathrm{D}}\mathbf{T}_{\mathrm{C}} = \left(^{\mathrm{W}}\mathbf{T}_{\mathrm{D}}\right)^{-1}\,{^{\mathrm{W}}\mathbf{T}_{\mathrm{C}}}.
\label{eq:extrinsik}
\end{equation}
Pro Kamera werden mehrere solcher Messungen gesammelt und gemittelt, um den Einfluss des Detektionsrauschens zu verringern.
Die Translation wird dabei arithmetisch gemittelt, die Orientierung nach dem in Abschnitt~\ref{sec:quaternionen_mittelung} beschriebenen Verfahren gemäß Gleichung~\eqref{eq:quaternion_mittelung}.
Das Ergebnis wird je Kamera als Transformation in einer YAML-Datei abgelegt.
Aus diesem Vorgehen ergeben sich folgende Vorteile:
Die Kameras können an einer beliebigen Stelle und in beliebiger Ausrichtung an der Drohne montiert werden, solange ein Tag gesehen wird. 
Wird eine Kamera getauscht oder neu an der Flugplattform positioniert, genügt ein neuer Kalibrierlauf, ohne die Kamera aufwändig neu einzumessen. 



\section{Positionsbestimmung durch AprilTags}
\label{sec:pos_positionsbestimmung}
Die Positionsbestimmung übernimmt der selbst entwickelte Knoten \texttt{vision\_pose\_uxrce\_node}.
Er berechnet aus den Markerdetektionen die Pose des Fluggeräts in der Karte und übergibt sie an die Kommunikation.
Der Knoten ist für den Betrieb mit zwei Kameras konzipiert, deren Detektionen unabhängig ausgewertet und anschließend zu einer 
gemeinsamen Schätzung fusioniert werden.
Der Ablauf gliedert sich in vier Schritte: die Posenberechnung je Marker, die Ausreißerfilterung, die gewichtete Fusion
und die Umrechnung auf den Drohnenmittelpunkt.
Das Verhalten des Knotens wird über eine Reihe von Startparametern gesteuert, die das Laden der Karte, die Fusion beider Kameras,
die Ausreißerfilterung und die Kovarianzschätzung konfigurieren.
Tabelle~\ref{tab:vision_pose_params} fasst diese Parameter thematisch geordnet zusammen.

{\footnotesize
\begin{longtable}{@{}l p{0.58\textwidth}@{}}
\caption{Startparameter des Knotens \texttt{vision\_pose\_uxrce\_node}}
\label{tab:vision_pose_params} \\
\toprule
\textbf{Parameter} & \textbf{Bedeutung} \\
\midrule
\endfirsthead

\multicolumn{2}{@{}l}{\footnotesize\itshape Fortsetzung von Tabelle~\ref{tab:vision_pose_params}} \\
\toprule
\textbf{Parameter} & \textbf{Bedeutung} \\
\midrule
\endhead

\midrule
\multicolumn{2}{r@{}}{\footnotesize\itshape Fortsetzung auf der nächsten Seite} \\
\endfoot

\bottomrule
\endlastfoot

\multicolumn{2}{@{}l}{\textit{Karte und Kalibrierung}} \\
\texttt{map\_file}                       & Pfad zur TagSLAM-Karte mit den Tag-Posen (Pflichtangabe) \\
\texttt{calibration\_file}               & Pfad zur YAML mit den Kamera-zu-Drohnenmittelpunkt-Transformationen (Pflichtangabe) \\
\addlinespace
\multicolumn{2}{@{}l}{\textit{Topics und Bezugssysteme}} \\
\texttt{cam1\_topic}, \texttt{cam2\_topic} & Eingangstopics der Tag-Detektionen beider Kameras \\
\texttt{cam1\_name}, \texttt{cam2\_name}   & Bezeichner der Kameras zur Zuordnung der Kalibrierung \\
\texttt{publish\_topic}                  & Ausgangstopic der \texttt{VehicleOdometry} an PX4 \\
\texttt{map\_frame}                      & Name des Karten-Koordinatensystems \\
\texttt{timesync\_topic}                 & Topic, auf dem PX4 den Timesync-Offset publiziert \\
\addlinespace
\multicolumn{2}{@{}l}{\textit{PX4-Anbindung und Zeit}} \\
\texttt{pose\_frame}                     & Referenzrahmen der gemeldeten Pose (1 = NED) \\
\texttt{apply\_timesync\_offset}         & Ob der PX4-Timesync-Offset auf die Zeitstempel angewendet wird \\
\addlinespace
\multicolumn{2}{@{}l}{\textit{Fusion und Glättung}} \\
\texttt{publish\_rate\_hz}               & Frequenz, mit der die fusionierte Pose publiziert wird \\
\texttt{max\_estimate\_age\_s}           & Maximalalter einer Kameraschätzung, bevor sie verworfen wird \\
\texttt{cam1\_weight}, \texttt{cam2\_weight} & Gewichtung der jeweiligen Kamera in der Fusion \\
\texttt{position\_smoothing\_alpha}      & Glättungsfaktor des EMA-Filters für die Position \\
\texttt{rotation\_smoothing\_alpha}      & Glättungsfaktor der Rotationsglättung (SLERP) \\
\texttt{min\_tags\_required}             & Mindestzahl sichtbarer Tags pro Frame zur Verwendung \\
\addlinespace
\multicolumn{2}{@{}l}{\textit{Ausreißerfilterung (RANSAC)}} \\
\texttt{ransac\_base\_threshold\_m}      & Inlier-Schwelle bei der Referenzdistanz \\
\texttt{ransac\_reference\_distance\_m}  & Distanz, bei der die Basisschwelle gilt \\
\texttt{ransac\_min\_inliers\_ratio}     & Mindestanteil konsistenter Tags für einen gültigen Konsens \\
\texttt{ransac\_max\_iterations}         & Maximale Iterationszahl (darunter deterministisch-vollständig) \\
\texttt{ransac\_refinement\_iterations}  & Nachverfeinerung der Schätzung mit den Inliern \\
\texttt{ransac\_max\_distance\_factor}   & Obergrenze des distanzabhängigen Skalierungsfaktors \\
\addlinespace
\multicolumn{2}{@{}l}{\textit{Kovarianzschätzung}} \\
\texttt{use\_covariance}                 & Ob geschätzte Varianzen in die Nachricht geschrieben werden \\
\texttt{covariance\_pos\_base\_sigma\_m} & Basis-Standardabweichung der Position unter Referenzbedingungen \\
\texttt{covariance\_rot\_base\_sigma\_rad} & Basis-Standardabweichung der Orientierung unter Referenzbedingungen \\
\texttt{covariance\_reference\_distance\_m} & Distanz, bei der die Basis-Standardabweichungen gelten \\
\texttt{covariance\_reference\_tag\_size\_m} & Tag-Größe, bei der die Basis-Standardabweichungen gelten \\
\texttt{covariance\_max\_distance\_factor} & Obergrenze des distanzabhängigen Skalierungsfaktors \\
\texttt{covariance\_max\_size\_factor}   & Obergrenze des größenabhängigen Skalierungsfaktors \\
\addlinespace
\multicolumn{2}{@{}l}{\textit{Multi-Size und Debug}} \\
\texttt{isaac\_reference\_size}          & Referenzgröße zur Translationskorrektur bei abweichenden Tag-Größen (0 = aus Karte) \\
\texttt{publish\_tag\_map\_tf}           & Ob die Tag-Karte als statische TF veröffentlicht wird (Debug) \\
\texttt{tag\_frame\_prefix}              & Präfix der TF-Framenamen der Tags \\
\end{longtable}
}





Hier fehlt noch, welche Dateien zum Start benötigt werden. 
%So fehlt map_file
%außerdem fehlt calibraiton_file
die ausrichtung des koordinatensystems fehlt auch noch. (z nach unten, oder z nach oben)

Im Folgenden werden die wichtigesten Opterationen nach und nach erläutert. 

\subsection{Posenberechnung je Marker}
\label{subsec:pos_posenberechnung}
Jeder im Kamerabild erkannte Marker liefert eine eigenständige Schätzung der Kamerapose im Kartenkoordinatensystem.
Die Größe und Lage eines Markers~$\mathrm{M}$ ist durch die beim Programmstart geladene Karte bekannt.
Die Lage wird im Kartenkoordinatensystem~$\mathrm{K}$ als $^{\mathrm{K}}\mathbf{T}_{\mathrm{M}}$ beschrieben.
Die Orientierung der Marker ist in der Karte als Rodrigues-Vektor hinterlegt und wird beim Laden gemäß
Gleichung~\eqref{eq:rodrigues_quaternion_eingesetzt} in ein Einheitsquaternion umgerechnet.
Wie in Abschnitt~\ref{sec:pos_detektion_apriltags} beschrieben, liefert die Detektion zu jedem Marker
dessen Lage relativ zum Kamerakoordinatensystem~$\mathrm{C}$ als Translationsvektor $^{\mathrm{C}}\tilde{\mathbf{t}}_{\mathrm{M}}$ und
Einheitsquaternion $^{\mathrm{C}}\mathbf{q}_{\mathrm{M}}$.

Im ersten Schritt wird aus dem detektierten Translationsvektor $^{\mathrm{C}}\tilde{\mathbf{t}}_{\mathrm{M}}$
der skalierte Vektor $^{\mathrm{C}}\mathbf{t}_{\mathrm{M}}$ berechnet, der für Marker, deren Größe von der
Referenz-Markergröße abweicht, entsprechend korrigiert wird.
Bezeichnet $a_{\mathrm{M}}$ die Kantenlänge des Markers und $a_{\mathrm{ref}}$ die Referenz-Markergröße, ergibt sich der skalierte Vektor zu
\begin{equation}
^{\mathrm{C}}\mathbf{t}_{\mathrm{M}} = \frac{a_{\mathrm{M}}}{a_{\mathrm{ref}}}\,{^{\mathrm{C}}\tilde{\mathbf{t}}_{\mathrm{M}}}.
\label{eq:pos_skalierung}
\end{equation}


%Im ersten Schritt werden die Translationsvektoren $^{\mathrm{C}}\mathbf{t}_{\mathrm{M}}$ der Marker, deren Größe von der
%Referenz-Markergröße abweicht, entsprechend skaliert.
%Bezeichnet $a_{\mathrm{M}}$ die Kantenlänge des Markers und $a_{\mathrm{ref}}$ die Referenz-Markergröße, ergibt sich der skalierte Vektor zu
%\begin{equation}
%^{\mathrm{C}}\mathbf{t}_{\mathrm{M}} \leftarrow \frac{a_{\mathrm{M}}}{a_{\mathrm{ref}}}\,{^{\mathrm{C}}\mathbf{t}_{\mathrm{M}}}.
%\label{eq:pos_skalierung}
%\end{equation}
Die Referenz-Markergröße $a_{\mathrm{ref}}$ wird über den Parameter \texttt{isaac\_reference\_size} übergeben.
Die Orientierung bleibt dabei unberührt, da die Größe des Markers keinen Einfluss darauf hat.
Der Isaac-Ros Detektor gibt die Orientierung in einer anderen Konvention aus als das in TagSlam verwendete \gls{ros2} standard paket.
Daher muss vor der weiteren Verarbeitung die detektierte Orientierung $^{\mathrm{C}}\mathbf{q}_{\mathrm{M}}$ durch eine feste
Drehung um $180^\circ$ um die $x$-Achse des Markers an die Konvention der Karte angeglichen werden.

Zur Positionsbestimmung  ist die Pose der Kamera im Kartenkoordinatensystem, $^{\mathrm{K}}\mathbf{T}_{\mathrm{C}}$ gesucht.
Konzeptionell entspricht dies der Verkettung der bekannten Markerpose mit der invertierten Detektion gemäß den
Gleichungen~\eqref{eq:trafo_verkettung} und~\eqref{eq:trafo_inverse}:
\begin{equation}
^{\mathrm{K}}\mathbf{T}_{\mathrm{C}} = {^{\mathrm{K}}\mathbf{T}_{\mathrm{M}}}\,\left(^{\mathrm{C}}\mathbf{T}_{\mathrm{M}}\right)^{-1}.
\label{eq:pos_kamera_in_karte}
\end{equation}
Hierfür wird zunächst die detektierte Pose invertiert, um die Lage der Kamera relativ zum Marker, $^{\mathrm{M}}\mathbf{T}_{\mathrm{C}}$, zu erhalten.
Die inverse Orientierung ergibt sich als Konjugation nach Gleichung~\eqref{eq:quaternion_konjungiert}, die inverse Translation durch Rotation
des negierten Vektors mit dieser Orientierung nach Gleichung~\eqref{eq:quaternion_rotation}:
\begin{align}
^{\mathrm{M}}\mathbf{q}_{\mathrm{C}} &= {^{\mathrm{C}}\bar{\mathbf{q}}_{\mathrm{M}}}, \label{eq:pos_inverse_q}\\
^{\mathrm{M}}\mathbf{t}_{\mathrm{C}} &= {^{\mathrm{C}}\bar{\mathbf{q}}_{\mathrm{M}}} \cdot \bigl(-{^{\mathrm{C}}\mathbf{t}_{\mathrm{M}}}\bigr) \cdot {^{\mathrm{C}}\mathbf{q}_{\mathrm{M}}}. \label{eq:pos_inverse_t}
\end{align}
Durch Verkettung dieser Pose mit der bekannten Markerpose folgt die gesuchte Kamerapose, wobei die Orientierungen nach Gleichung~\eqref{eq:quaternion_multiplikation} multipliziert und die Translation analog zu Gleichung~\eqref{eq:trafo_punkt} zusammengesetzt werden:
\begin{align}
^{\mathrm{K}}\mathbf{q}_{\mathrm{C}} &= {^{\mathrm{K}}\mathbf{q}_{\mathrm{M}}} \cdot {^{\mathrm{M}}\mathbf{q}_{\mathrm{C}}}, \label{eq:pos_kompose_q}\\
^{\mathrm{K}}\mathbf{t}_{\mathrm{C}} &= {^{\mathrm{K}}\mathbf{t}_{\mathrm{M}}} + {^{\mathrm{K}}\mathbf{q}_{\mathrm{M}}} \cdot {^{\mathrm{M}}\mathbf{t}_{\mathrm{C}}} \cdot {^{\mathrm{K}}\bar{\mathbf{q}}_{\mathrm{M}}}. \label{eq:pos_kompose_t}
\end{align}
Jeder sichtbare Marker liefert auf diese Weise eine vollständige Schätzung der Kamerapose $\bigl(^{\mathrm{K}}\mathbf{t}_{\mathrm{C}},\,{^{\mathrm{K}}\mathbf{q}_{\mathrm{C}}}\bigr)$.
%Da nahe Marker eine genauere Schätzung erlauben als weiter entfernte, wird jeder Schätzung ein Gewicht zugeordnet, das mit dem Quadrat der Distanz~$d$ zwischen Kamera und Marker abnimmt:
%\begin{equation}
%d = \lVert{^{\mathrm{C}}\mathbf{t}_{\mathrm{M}}}\rVert, \qquad w = \frac{1}{d^{2}}.
%\label{eq:pos_gewicht}
%\end{equation}
Marker, die im Bild größer erscheinen, erlauben eine genauere Schätzung als kleiner erscheinende, da sie auf mehr Pixel abgebildet werden.
Jeder Schätzung wird daher ein Gewicht zugeordnet, das mit der scheinbaren Größe des Markers im Bild zunimmt.
Als Maß hierfür dient die Norm der unmittelbar detektierten, \emph{unskalierten} Translation $^{\mathrm{C}}\tilde{\mathbf{t}}_{\mathrm{M}}$:
\begin{equation}
\tilde{d} = \lVert{^{\mathrm{C}}\tilde{\mathbf{t}}_{\mathrm{M}}}\rVert, \qquad w = \frac{1}{\tilde{d}^{2}}.
\label{eq:pos_gewicht}
\end{equation}
Da der Detektor die Translation für alle Marker mit der festen Referenzgröße $a_{\mathrm{ref}}$ berechnet, 
ist die Norm $\tilde{d}$ kein Maß für die tatsächliche Entfernung, sondern für die scheinbare Größe des Markers im Bild.
So wird abgebildet, dass im Bild große Marker aufgrund der höheren Genauigkeit mit einem größeren Gewicht in die Schätzung einfließt, 
als kleinere Marker mit geringerer Genauigkeit.
Die Gewichtung der Marker erfolgt dadurch nur auf Basis der Pixelgröße, unabhängig von der tatsächlichen Entfernung des markers zum Objektiv.



\subsection{Ausreißerfilterung}
\label{subsec:pos_ausreisserfilterung}
Aus dem vorigen Schritt liegt für jedes Bild und jeden darin sichtbaren Marker eine eigenständige Schätzung
der Kameraposition $^{\mathrm{K}}\mathbf{t}_{\mathrm{C}}$ im Kartenkoordinatensystem vor.
Bei korrekten Detektionen stimmen diese Schätzungen überein, da sie dieselbe physische Kamera beschreiben.
Eine fehlerhafte Detektion oder eine ungenaue Karte führt dagegen zu einer abweichenden Schätzung.
Zur Bestimmung der größten konsistenten Teilmenge wird das in Abschnitt~\ref{sec:ransac} beschriebene \gls{ransac}-Verfahren angewendet,
das ausschließlich die Translationen vergleicht.
Da die Translation nach Gleichung~\eqref{eq:pos_kompose_t} aus der Rotation berechnet wird, schlägt sich eine fehlerhafte Rotation
unmittelbar in einer fehlerhaften Translation nieder, sodass der Vergleich der Translationen genügt, um beide Fehlerarten zu erfassen.

Die minimale Teilmenge zur Modellbestimmung der Position besteht aus einer einzelnen Schätzung, die als Hypothese für die Kameraposition dient.
Als Abstand $d(\mathbf{x}_i, \mathcal{M})$ nach Gleichung~\eqref{eq:ransac_inlier} dient die euklidische Distanz zwischen einer Schätzung und der Hypothese.
Da das \gls{uav} sehr unterschiedlich weit von Tags entfernt sein kann und mit der Entfernung die Positioniergenauigkeit in Metern abnimmt,
wird der Inlier-Schwellwert variabel zur mittleren Markerdistanz gestaltet.
Bezeichnet $\bar{d}$ die mittlere Markerdistanz, $d_{\mathrm{ref}}$ die Referenzdistanz (Parameter \texttt{ransac\_reference\_\-distance\_m})
und $\tau_0$ die Basisschwelle (Parameter \texttt{ransac\_base\_\-threshold\_m}), ergibt sich der adaptive Schwellenwert $\tau$ zu
\begin{equation}
\tau = \tau_0 \cdot f, \qquad
f =
\begin{cases}
0{,}5 & \text{falls } \left(\dfrac{\bar{d}}{d_{\mathrm{ref}}}\right)^{2} < 0{,}5, \\[2ex]
f_{\max} & \text{falls } \left(\dfrac{\bar{d}}{d_{\mathrm{ref}}}\right)^{2} > f_{\max}, \\[2ex]
\left(\dfrac{\bar{d}}{d_{\mathrm{ref}}}\right)^{2} & \text{sonst.}
\end{cases}
\label{eq:pos_ransac_schwelle}
\end{equation}
Die untere Schranke verhindert einen unrealistisch kleinen Schwellenwert bei sehr nahen Markern, 
die obere Schranke $f_{\max}$ (Parameter \texttt{ransac\_max\_distance\_factor}) begrenzt ihn bei großer Distanz.
Die weitere Verarbeitung hängt von der Anzahl~$n$ der vorliegenden Schätzungen ab.

Wird in einem Bild nur ein einzelner Marker detektiert, ist kein vergleich möglich. 
Die Schätzung wird unverändert weitergegeben.

Bei genau zwei Schätzungen wird geprüft, ob ihr Abstand den Schwellenwert~$\tau$ unterschreitet. 
Ist dies der Fall, werden beide übernommen, andernfalls wird keine Konsensmenge ausgegeben und die Schätzungen verworfen.

Liegen mindestens drei Schätzungen vor, wird in jeder Iteration eine Schätzung als Hypothese herangezogen und gezählt, 
wie viele der übrigen Schätzungen nach dem Kriterium aus Gleichung~\eqref{eq:ransac_inlier} innerhalb des Schwellenwerts $\tau$ liegen.
Die Anzahl der Iterationen hängt von der Schätzungsanzahl~$n$ und dem Parameter \texttt{ransac\_max\_iterations} ab.
Die Hypothese mit der größten Konsensmenge bestimmt die Inlier-Menge.

Die Konsensmenge wird nur dann ausgegeben, wenn sie als ausreichend vertrauenswürdig gilt.
Die erforderliche Mindestanzahl $n_{\min}$ wird durch zwei Parameter festgelegt: die 
Mindestmarkeranzahl $n_{\mathrm{Tags}}$ (Parameter \texttt{min\_tags\_required}) und einen Anteil $r$ aller 
Schätzungen (Parameter \texttt{ransac\_min\_inliers\_ratio}).
Maßgeblich ist der größere der beiden Werte:
\begin{equation}
n_{\min} = \max\!\left(n_{\mathrm{Tags}},\; \lceil r \cdot n \rceil\right).
\label{eq:pos_ransac_konsens}
\end{equation}
Wird $n_{\min}$ nicht erreicht, gilt der Konsens als unzuverlässig und es wird keine Konsensmenge ausgegeben.



\subsection{Gewichtete Fusion der Konsensmenge}
\label{subsec:pos_fusion}
Die aus der Ausreißerfilterung hervorgegangene Konsensmenge $\mathcal{I}$ enthält die als zuverlässig eingestuften Schätzungen der Kamerapose.
Diese werden zu einer einzigen Schätzung $\bigl(^{\mathrm{K}}\mathbf{t}_{\mathrm{C}},\,{^{\mathrm{K}}\mathbf{q}_{\mathrm{C}}}\bigr)$ zusammengefasst.
Da sich Position und Orientierung unterschiedlich mitteln lassen, werden beide getrennt behandelt.
In beide Fälle gehen die Schätzungen mit den Gewichten $w_i$ aus Gleichung~\eqref{eq:pos_gewicht} ein, sodass im Bild größer erscheinende Marker stärker zur fusionierten Pose beitragen.

Die Position ergibt sich als gewichteter Mittelwert der einzelnen Positionsschätzungen $^{\mathrm{K}}\mathbf{t}_{\mathrm{C}}^{(i)}$:
\begin{equation}
^{\mathrm{K}}\mathbf{t}_{\mathrm{C}} = \frac{\displaystyle\sum_{i \in \mathcal{I}} w_i \, {^{\mathrm{K}}\mathbf{t}_{\mathrm{C}}^{(i)}}}{\displaystyle\sum_{i \in \mathcal{I}} w_i}.
\label{eq:pos_fusion_t}
\end{equation}

Die Orientierung wird nach dem in Abschnitt~\ref{sec:quaternionen_mittelung} beschriebenen Verfahren zur Mittelung von Quaternionen zusammengefasst.
Als Referenzquaternion $\mathbf{q}_{\mathrm{ref}}$ für Gleichung~\eqref{eq:quaternion_vorzeichen} dient die erste Schätzung der Konsensmenge.
\begin{equation}
^{\mathrm{K}}\mathbf{q}_{\mathrm{C}} = \frac{\displaystyle\sum_{i \in \mathcal{I}} w_i \, {^{\mathrm{K}}\mathbf{q}_{\mathrm{C}}^{(i)}}}{\left\lVert \displaystyle\sum_{i \in \mathcal{I}} w_i \, {^{\mathrm{K}}\mathbf{q}_{\mathrm{C}}^{(i)}} \right\rVert}.
\label{eq:pos_fusion_q}
\end{equation}

Damit liegt für das jeweilige Kamerabild eine fusionierte Schätzung der 
Kamerapose $\bigl(^{\mathrm{K}}\mathbf{t}_{\mathrm{C}},\,{^{\mathrm{K}}\mathbf{q}_{\mathrm{C}}}\bigr)$ im Kartenkoordinatensystem vor.




\subsection{Umrechnung der Fusion auf den Drohnenmittelpunkt}
\label{subsec:pos_drohnenmittelpunkt}
Die fusionierte Pose aus den Gleichungen~\eqref{eq:pos_fusion_t} und~\eqref{eq:pos_fusion_q} beschreibt die Lage der Kamera im Kartenkoordinatensystem.
Für die Regelung des \gls{uav} ist jedoch die Pose des Drohnenmittelpunkts~$\mathrm{D}$ maßgeblich.
Da die Kamera starr an der Drohne montiert ist, ist die Transformation $^{\mathrm{C}}\mathbf{T}_{\mathrm{D}}$ 
von der jeweiligen Kamera zum Drohnenmittelpunkt fest und aus der Kalibrierung bekannt.
Die gesuchte Pose $^{\mathrm{K}}\mathbf{T}_{\mathrm{D}}$ ergibt sich durch Verkettung der fusionierten 
Kamerapose mit dieser Transformation gemäß Gleichung~\eqref{eq:trafo_verkettung}:
\begin{equation}
^{\mathrm{K}}\mathbf{T}_{\mathrm{D}} = {^{\mathrm{K}}\mathbf{T}_{\mathrm{C}}}\,{^{\mathrm{C}}\mathbf{T}_{\mathrm{D}}}.
\label{eq:pos_drohne}
\end{equation}
Die Orientierung wird dabei nach Gleichung~\eqref{eq:quaternion_multiplikation} multipliziert und die Translation analog zu Gleichung~\eqref{eq:trafo_punkt} zusammengesetzt:
\begin{align}
^{\mathrm{K}}\mathbf{q}_{\mathrm{D}} &= {^{\mathrm{K}}\mathbf{q}_{\mathrm{C}}} \cdot {^{\mathrm{C}}\mathbf{q}_{\mathrm{D}}}, \label{eq:pos_drohne_q}\\
^{\mathrm{K}}\mathbf{t}_{\mathrm{D}} &= {^{\mathrm{K}}\mathbf{t}_{\mathrm{C}}} + {^{\mathrm{K}}\mathbf{q}_{\mathrm{C}}} \cdot {^{\mathrm{C}}\mathbf{t}_{\mathrm{D}}} \cdot {^{\mathrm{K}}\bar{\mathbf{q}}_{\mathrm{C}}}.
\label{eq:pos_drohne_t}
\end{align}
Damit liegt für jede Kamera eine Schätzung der Pose des Drohnenmittelpunkts im Kartenkoordinatensystem vor.


\subsection{Fusion der Kameras}
\label{subsec:pos_kamerafusion}
Die bisher beschriebenen Schritte der Positionierung -- Detektion, Skalierung, Posenberechnung je Marker, Ausreißerfilterung, 
gewichtete Fusion der Konsensmenge und Umrechnung auf den Drohnenmittelpunkt -- werden für jede Kamera und für jedes Kamerabild durchlaufen.
Jede Kamera liefert dadurch fortlaufend eine eigene Schätzung der Pose des Drohnenmittelpunkts $^{\mathrm{K}}\mathbf{T}_{\mathrm{D}}$ 
nach Gleichung~\eqref{eq:pos_drohne} im Kartenkoordinatensystem.
Da die Kameras unabhängig voneinander und mit unterschiedlicher Bildrate arbeiten können, treffen diese Schätzungen asynchron ein.

Die Zusammenführung erfolgt daher nicht bei jeder eingehenden Detektion, sondern in festem Takt mit der Frequenz~$\mathtt{publish\_rate\_hz}$.
Bei jedem Aufruf wird je Kamera die jeweils aktuellste Schätzung herangezogen.
Eine Schätzung wird nur berücksichtigt, wenn ihr Alter den Parameter \texttt{max\_estimate\_age\_s} nicht überschreitet, veraltete Schätzungen werden verworfen.
Liegt keine gültige Schätzung vor, wird keine Pose ausgegeben.

Den beiden Kameras ist jeweils ein fester Gewichtungsfaktor $w_1$ bzw. $w_2$ zugeordnet (Parameter \texttt{cam1\_weight} bzw. \texttt{cam2\_weight}), über den eine Kamera als zuverlässiger eingestuft werden kann.
Anders als das von der scheinbaren Markergröße abhängige Gewicht aus Gleichung~\eqref{eq:pos_gewicht} sind diese beiden Faktoren fest vorgegeben.
Die fusionierte Position ergibt sich analog zu Gleichung~\eqref{eq:pos_fusion_t} als gewichteter Mittelwert der beiden Schätzungen $^{\mathrm{K}}\mathbf{t}_{\mathrm{D}}^{(1)}$ und $^{\mathrm{K}}\mathbf{t}_{\mathrm{D}}^{(2)}$:
\begin{equation}
^{\mathrm{K}}\mathbf{t}_{\mathrm{D}} = \frac{w_1\,{^{\mathrm{K}}\mathbf{t}_{\mathrm{D}}^{(1)}} + w_2\,{^{\mathrm{K}}\mathbf{t}_{\mathrm{D}}^{(2)}}}{w_1 + w_2}.
\label{eq:pos_kamerafusion_t}
\end{equation}
Die Orientierung wird nach dem in Abschnitt~\ref{sec:quaternionen_mittelung} beschriebenen Verfahren zusammengefasst.
Als Referenzquaternion dient die Schätzung der ersten Kamera $^{\mathrm{K}}\mathbf{q}_{\mathrm{D}}^{(1)}$.
Das Vorzeichen von $^{\mathrm{K}}\mathbf{q}_{\mathrm{D}}^{(2)}$ wird nach Gleichung~\eqref{eq:quaternion_vorzeichen} an dieses angeglichen und beide anschließend gemäß Gleichung~\eqref{eq:quaternion_mittelung} mit den Gewichten $w_1$ und $w_2$ gemittelt:
\begin{equation}
^{\mathrm{K}}\mathbf{q}_{\mathrm{D}} = \frac{w_1\,{^{\mathrm{K}}\mathbf{q}_{\mathrm{D}}^{(1)}} + w_2\,{^{\mathrm{K}}\mathbf{q}_{\mathrm{D}}^{(2)}}}{\left\lVert w_1\,{^{\mathrm{K}}\mathbf{q}_{\mathrm{D}}^{(1)}} + w_2\,{^{\mathrm{K}}\mathbf{q}_{\mathrm{D}}^{(2)}} \right\rVert}.
\label{eq:pos_kamerafusion_q}
\end{equation}
Liefert nur eine der beiden Kameras eine gültige Schätzung, wird deren Pose unverändert übernommen.


\subsection{Zeitliche Glättung der fusionierten Pose}
\label{subsec:pos_glaettung}
Die nach Gleichung~\eqref{eq:pos_kamerafusion_t} und~\eqref{eq:pos_kamerafusion_q} fusionierte Pose wird bei jedem 
Taktzeitpunkt aus den jeweils aktuellsten Kameraschätzungen neu bestimmt.
Dadurch übertragen sich das Rauschen der einzelnen Detektionen und kleine Sprünge zwischen aufeinanderfolgenden Schätzungen unmittelbar auf die ausgegebene Pose.
Um diese hochfrequenten Anteile zu dämpfen, wird die fusionierte Pose vor der Ausgabe zeitlich geglättet.
Die Gelättung von Position und Orientierung getrennt. 

Die Position wird durch einen exponentiell gewichteten gleitenden Mittelwert, \gls{ema} geglättet.
Bezeichnet $^{\mathrm{K}}\mathbf{t}_{\mathrm{D}}[k]$ die fusionierte Position zum Taktzeitpunkt~$k$ und $^{\mathrm{K}}\hat{\mathbf{t}}_{\mathrm{D}}[k]$ die geglättete Position, so ergibt sich diese rekursiv aus dem vorherigen geglätteten Wert und der aktuellen Schätzung:
\begin{equation}
^{\mathrm{K}}\hat{\mathbf{t}}_{\mathrm{D}}[k] = {^{\mathrm{K}}\hat{\mathbf{t}}_{\mathrm{D}}[k-1]} + \alpha_{\mathrm{t}}\left({^{\mathrm{K}}\mathbf{t}_{\mathrm{D}}[k]} - {^{\mathrm{K}}\hat{\mathbf{t}}_{\mathrm{D}}[k-1]}\right).
\label{eq:pos_ema}
\end{equation}
Der Glättungsfaktor $\alpha_{\mathrm{t}} \in [0,1]$ wird über den Parameter \texttt{position\_smoothing\_alpha} festgelegt.
Ein Wert von $\alpha_{\mathrm{t}} = 1$ entspricht keiner Glättung, kleinere Werte dämpfen die Schätzung stärker und erhöhen zugleich deren Trägheit.

Die Orientierung wird durch \gls{slerp} nach Gleichung~\eqref{eq:slerp} zwischen dem vorherigen geglätteten Quaternion $^{\mathrm{K}}\hat{\mathbf{q}}_{\mathrm{D}}[k-1]$ 
und der aktuellen Schätzung $^{\mathrm{K}}\mathbf{q}_{\mathrm{D}}[k]$ geglättet:
\begin{equation}
^{\mathrm{K}}\hat{\mathbf{q}}_{\mathrm{D}}[k] = \frac{\sin\!\bigl((1-\alpha_{\mathrm{q}})\,\theta\bigr)}{\sin\theta}\,{^{\mathrm{K}}\hat{\mathbf{q}}_{\mathrm{D}}[k-1]} + \frac{\sin(\alpha_{\mathrm{q}}\,\theta)}{\sin\theta}\,{^{\mathrm{K}}\mathbf{q}_{\mathrm{D}}[k]}.
\label{eq:pos_slerp}
\end{equation}
Dabei ist $\theta$ der Winkel zwischen dem vorherigen geglätteten Quaternion und der aktuellen Schätzung.
Der Interpolationsfaktor $\alpha_{\mathrm{q}} \in [0,1]$ wird über den Parameter \texttt{rotation\_smoothing\_alpha} 
festgelegt und wirkt analog zu $\alpha_{\mathrm{t}}$.


\subsection{Umrechnung in das PX4-Koordinatensystem und Ausgabe}
\label{subsec:pos_frame_ausgabe}
Die bisher bestimmte geglättete Pose $\bigl(^{\mathrm{K}}\hat{\mathbf{t}}_{\mathrm{D}},\,{^{\mathrm{K}}\hat{\mathbf{q}}_{\mathrm{D}}}\bigr)$ 
liegt im Kartenkoordinatensystem vor, das der in Abschnitt~\ref{sec:koordinatensysteme} eingeführten ENU/FLU-Konvention folgt.
Da PX4 die Pose im NED/FRD-System benötigt, wird sie vor der Ausgabe von ENU/FLU nach NED/FRD umgerechnet.
Anschließend wird die umgerechnete Pose über die uXRCE-DDS-Schnittstelle auf dem Topic \topic{/fmu/in/vehicle\_visual\_odometry} an PX4 übergeben.

\section{hier muss noch korrektur gelesen werden!!!}
\section{Schnittstelle zwischen ROS2 und PX4}
\subsection{Kommunikation zwischen ROS2 und PX4}
\label{subsec:kommunikation_ros2_px4}
Die Kommunikation zwischen ROS2 und \gls{px4} erfolgt über die uXRCE-DDS-Schnittstelle.
Auf \gls{px4} läuft dabei der uXRCE-DDS-Client, auf dem Companion-Computer der zugehörige Agent.
Der Agent wird über das Launch-File als Prozess \texttt{MicroXRCEAgent} gestartet und kommuniziert per UDP über die Ethernet-Verbindung auf einem konfigurierbaren Port.
Über diese Verbindung werden die uORB-Topics von \gls{px4} als ROS2-Topics bereitgestellt und umgekehrt.
So kann der Positionierungs-Node die geschätzte Pose als \texttt{VehicleOdometry} auf dem Topic \topic{/fmu/in/vehicle\_visual\_odometry} an \gls{px4} übergeben.

